\documentclass[11pt]{article}

\usepackage[a4paper,margin=1in]{geometry}
\usepackage[T1]{fontenc}
\usepackage{lmodern}
\usepackage{microtype}
\usepackage{parskip}

\setlength{\emergencystretch}{2em}

\title{\vspace{-1.5em}Physics-based protein--ligand docking at scale:\\
RosettaLigand samples near-native poses but cannot reliably rank them}
\author{%
  Paul Eisenhuth\textsuperscript{1,$\ast$}, Luca Friebe\textsuperscript{1,$\ast$},
  Josephine Eboh\textsuperscript{1},
  Chris Jurich\textsuperscript{2}, Rocco Moretti\textsuperscript{2},
  Jens Meiler\textsuperscript{1,2}\\[0.7em]
  {\small\textsuperscript{1}Leipzig University \qquad
   \textsuperscript{2}Vanderbilt University}\\[0.3em]
  {\small\textsuperscript{$\ast$}These authors contributed equally.}}
\date{}

\begin{document}
\maketitle
\vspace{-2.5em}

\begin{center}
\large\textbf{Abstract}
\end{center}

\noindent
RosettaLigand is a popular physics-based tool for predicting how a drug-like
molecule binds to a protein. Many researchers rely on it, but no one has tested
carefully how well it actually works across large, modern benchmark sets. We ran
the largest such test so far, re-docking thousands of protein--ligand complexes
from three standard datasets (PDBbind, Astex Diverse, and PoseBusters) and
recording every predicted pose along with its energy scores. We find that
RosettaLigand is very good at \emph{finding} correct binding poses---for most
complexes, a near-correct pose is among its predictions, often beating recent
deep-learning methods. The problem is \emph{picking} the right one: its scoring
function frequently rates a wrong pose higher than the correct one. We tried to
fix this by re-weighting the energy terms and training machine-learning models to
score poses, but the improvements were small. A closer look at the energy
landscape suggests why: there are many different poses with almost identical
energy scores, so the current scoring cannot reliably tell them apart. We also
examine how well the scores predict actual binding strength, and we release a
large dataset of scored, energy-minimised structures for others to build on.
\emph{(Work in progress.)}

\end{document}
